\subsection{代数扩张}

定义1. —— 域K的扩域E称为（在K上）代数的，如果E的每个元素都在K上是代数的。域K的扩域E若不是代数的，则称为（在K上）超越的。

命题1. --- 要使K的扩域E成为代数扩张，其充分必要条件是E的每一个子K-代数A都是一个域。

条件是必要的：如果 E 在 K 上是代数的，且 \(x \neq 0\) 是 E 的子 K-代数 A 的一个元素，则由 V，第 16 页，定理 1，d) 可知 \(x^{-1} \in K[x] \subset A\) 。因此 A 是一个域。

条件是充分的：如果它被满足且 x 是 E 中一个 \(\neq 0\) 的元素，则环 K{[}x{]} 是一个域，因此 \(x^{-1} \in K[x]\) ；换言之，存在一个多项式 \(g \in K[X]\) 使得 \(x^{-1} = g(x)\) ，或者也有 \(xg(x) - 1 = 0\) ；这表明 x 在 K 上是代数的，因此 E 是 K 的一个代数扩张。

命题2. --- 如果K的一个扩张E的次数为有限数n，则E是代数的，且E中每个元素在K上的次数整除n。

因为，给定 \(x \in \mathbb{E}\) , \([\mathbf{K}(x): \mathbf{K}]\) 是有限的且整除 \(n(\mathbf{V}, p. 10, \text{Cor. } 1)\) ，因此 \(x\) 在 \(\mathbf{K}(\mathbf{V}, p. 17, \text{Cor. } 1)\) 上是代数的。

* 存在无限次数的代数扩张，例如有限域的代数闭包（V，第24页，注记4）。*

定理2. --- 设E是K的有限生成扩张，由元素 \(a_1, \ldots, a_m\) 生成，这些元素是K上的代数元；则E是K的有限次扩张。若 \(a_i\) 在 K 上（相对于 \(a_1, a_2, \ldots, a_{i-1}\) ) 的次数是 \(n_i\) （对于 \(1 \leq i \leq m\) )，则 E 在 K 上的次数为 \(n_1 n_2 \ldots n_m\) 且元素 \(a_1^{\nu_1} a_2^{\nu_2} \ldots a_m^{\nu_m}\) ( \(0 \leq \nu_i \leq n_i - 1\) ) 构成 E 在 K 上的一组基。

元素 \(a_i^{\nu_i}\) ( \(0 \leqslant \nu_i \leqslant n_i - 1\) ) 构成 \(K(a_1, a_2, \ldots, a_i)\) 在 \(K(a_1, a_2, \ldots, a_{i-1})\) 上的一组基，这由 V 第 16 页定理 1 b) 可知；因此该定理通过对 \(m\) 归纳，由第二卷第 222 页命题 25 得出。

推论1. --- 设E是K的一个扩域，且A是E的一个由在K上代数的元素组成的子集。则K(A)在K上是代数的，并且我们有K{[}A{]} = K(A)。

因为每个 \(x \in \mathrm{K}(\mathrm{A})\) 都属于一个域 \(\mathrm{K}(\mathrm{F})\) ，其中 F 是 A 的有限子集（V，第 11 页，推论）；现在 \(\mathrm{K}(\mathrm{F})\) 在 K 上是代数的，且由定理2等于 \(\mathrm{K}[\mathrm{F}]\) ，因此x在K上是代数的，且 \(\mathrm{K}(\mathrm{A}) = \mathrm{K}[\mathrm{A}]\) 。

推论2. --- 设L是K的一个扩张，E、F是L的子扩张。若F在K上是代数的，则子环K{[}E, F{]} 由 E ∪ F 生成的 L 是一个域，它等于 E(F)，并且在 E 上是代数的。

F 的每个元素在 K 上是代数的，因此在 E 上也是代数的（V，第 17 页，推论 2），从而 \(E(F)\) 是 E 的一个代数扩张，并且我们有 \(E(F) = E[F]\) 由推论1可知。

注记。--- 1) 采用定理 2 的记号，E = K{[}a₁, a₂, \ldots, aₘ{]} 因此 E 同构于商 K{[}X₁, X₂, \ldots, Xₘ{]}/α；由于 E 是域，α 是 K{[}X₁, \ldots, Xₘ{]} 中的极大理想。

\begin{enumerate}
\def\labelenumi{\arabic{enumi})}
\setcounter{enumi}{1}
\tightlist
\item
  设 \(\mathbf{E}\) 是 \(\mathbf{K}\) 的代数扩张，且次数无限。由定理2，存在一个无限序列 \((a_n)_{n \geq 1}\) ，其元素属于 \(\mathbf{E}\) ，使得 \(a_n \notin \mathbf{K}(a_1, a_2, \ldots, a_{n-1})\) ；定理2进一步表明 \(\mathbf{K}(a_1, a_2, \ldots, a_n)\) 在 \(\mathbf{K}\) 上的次数取任意大的值。换句话说，如果E是K的代数扩张，且次数 \([F:K]\) （即 E 的在 K 上为有限次的子扩张 F 的次数）有界，则 E 是 K 的有限次扩张。
\end{enumerate}

